% =====================================================================
%  DRAFT MANUSCRIPT v0.3 -- INCOMPLETE. NOT FOR CIRCULATION.
%  Autonomous Driving Agent / CARLA 0.9.16
%  Compile on Overleaf with pdfLaTeX. Class: IEEEtran (conference).
%
%  STATUS OF THIS FILE
%    Sections I-V   : drafted
%    Section  VI    : RENDERING FAILURES -- drafted, needs measured numbers
%    Section  VII   : RESULTS -- deliberately empty, corpus does not exist
%    Section  VIII  : DISCUSSION -- stub
%    Section  IX    : CONCLUSION -- unwritten
%    References     : partial, renumber after additions
%  Search for "TODO" to find every open item.
% =====================================================================

\documentclass[conference]{IEEEtran}
\IEEEoverridecommandlockouts

\usepackage{cite}
\usepackage{amsmath,amssymb}
\usepackage{graphicx}
\usepackage{booktabs}
\usepackage{array}
\newcolumntype{L}[1]{>{\raggedright\arraybackslash}p{#1}}
\usepackage{xcolor}
\usepackage{textcomp}
\usepackage{framed}
\usepackage[hidelinks]{hyperref}

% --- helper for the visible "unfinished" markers -------------------
% Open items appear as grey side-barred notes. Delete the \TODO calls
% (and this definition) once each section is finished.
\colorlet{shadecolor}{white}
\newcommand{\TODO}[1]{%
  \begingroup
  \setlength{\FrameSep}{4pt}%
  \def\FrameCommand{{\color{gray!55}\vrule width 1.2pt}\hspace{5pt}}%
  \MakeFramed{\advance\hsize-\width \FrameRestore}%
  \noindent\scriptsize\itshape\color{gray!80}#1%
  \endMakeFramed
  \endgroup}

\newcommand{\UNWRITTEN}[1]{\par\noindent{\itshape\color{gray!85}#1}\par\vspace{2pt}}

% Same kind of note, but breakable across a column (a framed box is not).
\newcommand{\TODOflow}[1]{%
  \par\noindent{\scriptsize\itshape\color{gray!80}#1\par}}

\begin{document}

\title{A Contract-Gated Imitation Learning Stack for Autonomous Driving
in CARLA, and the Rendering Failures That Gate Its Data Collection\\[3pt]
{\normalsize\color{red!65!black}\bfseries DRAFT v0.3 --- INCOMPLETE MANUSCRIPT --- NOT FOR CIRCULATION}}

\author{
\IEEEauthorblockN{
Medhansh Khattar (25BAI10672),\quad
Prashant Kumar Jaglan (25BAI10750),\quad
Yash Yadav (25BAI10545)\\[2pt]
Samyak Jain (25BAI10746),\quad
Sanchit Kumar Sinha (25BAI10229),\quad
Aryan Garg (25BAI10215)}
\IEEEauthorblockA{School of Computer Science and Engineering\\
Vellore Institute of Technology, Bhopal, India}
}

\maketitle

% =====================================================================
\begin{abstract}
Learning a driving policy from demonstration is usually reported as a
modelling result. In our experience the binding constraint has been
elsewhere: we cannot reliably render the environment. This paper
describes an imitation-learning stack built on CARLA 0.9.16 around a
deliberately small state-based policy (an
$8\rightarrow24\rightarrow16\rightarrow3$ perceptron, 667 parameters)
and three explicit contracts: an eight-field observation contract, a
seven-field action protocol, and a deployability contract an artifact
must satisfy before it may drive. An independent safety supervisor,
outside the learned policy, retains full braking authority. We report
the demonstration protocol and the route-level split that avoids
frame-level leakage; we then report, at length, the rendering
failures that dominate the project. The final camera image
cannot be rendered at speed on an integrated GPU, environment streaming
stalls the simulator, and the resulting crashes destroy episodes
mid-route. Because our acceptance rule discards any interrupted route,
these crashes propagate directly into an empty training corpus.
\textit{Quantitative evaluation is therefore unavailable: the accepted
corpus stands at zero routes of a twelve-route target, and no artifact
satisfies the v2 contract. Results are reserved for Sections~VII--IX and
remain unwritten.}
\end{abstract}

\begin{IEEEkeywords}
autonomous driving, imitation learning, behaviour cloning, CARLA,
rendering bottleneck, safety supervision, reproducibility.
\end{IEEEkeywords}

% =====================================================================
\section{Introduction}

Simulation is the default proving ground for learned driving policies,
and CARLA~\cite{carla} is the most widely adopted open platform: it
exposes an Unreal Engine world over an RPC interface, supplies OpenDRIVE
road topology, traffic and weather, and permits deterministic replay of
a seeded route. Failures cost nothing --- provided the simulation runs
at all.

Much of the literature following Bojarski \textit{et al.}~\cite{bojarski}
maps camera pixels directly to control. We take a narrower first step:
our policy consumes eight route-relative scalars rather than images. The
motivation is the size of the demonstration corpus a six-person
undergraduate team can realistically curate --- with a few thousand
frames a larger network memorises the road geometry of the training town
instead of learning to hold a lane~\cite{limits}. Our contribution is an
engineering one: the architecture, the contracts holding it together,
and --- centrally --- the rendering and stability failures met in
building it. Where a driving result would ordinarily appear, this draft
states plainly that it does not yet exist, and Section~VI explains why.

% =====================================================================
\section{Related Work}

End-to-end behaviour cloning for driving was demonstrated at scale
by~\cite{bojarski}. Codevilla \textit{et al.}~\cite{cil} added a
high-level command input for intersection ambiguity, and later
documented behaviour cloning's own limitations~\cite{limits}:
distributional shift and bias toward the modal action. Ross \textit{et
al.}~\cite{dagger} give the standard remedy, DAgger. Within CARLA,
Learning by Cheating~\cite{lbc} shows privileged-to-sensorimotor
distillation outperforming direct cloning, and
TransFuser~\cite{transfuser} sets the multi-modal baseline; our
observation contract is closer to the privileged stage of~\cite{lbc}.

\TODO{TODO --- RELATED WORK. Not yet cited, still to be read and added:
(i) Toromanoff \textit{et al.}, implicit affordances (CVPR 2020), and
where our reward logging sits relative to it; (ii) Kiran \textit{et
al.}, DRL survey (T-ITS 2021), for positioning; (iii) SAE J3016, so the
autonomy claim is precise; (iv) prior work on simulator throughput and
headless rendering --- unsurveyed, and directly relevant to Section~VI.
Owner: unassigned.}

% =====================================================================
\section{System Architecture}

The stack spans two machines on a direct Ethernet link. Machine~1 runs
the CARLA server and owns the world: physics, the OpenDRIVE map, traffic
manager, sensors and the vehicle. Machine~2 runs policy inference, the
route planner, an observer window and network visualisation. Each
iteration: machine~1 publishes world state; machine~2 runs the planner
and policy; machine~2 returns a driving command; machine~1's safety
layer validates it and only then applies it.

The map is deliberately split: CARLA's physical map stays on machine~1
because the simulator owns the world, while the navigation map that
picks a destination and computes waypoints may live on machine~2,
needing only road topology. A watchdog on machine~1 should brake fully
if the link drops; it is implemented but untested under induced failure.

% =====================================================================
\section{Methodology}

\subsection{Observation contract}

One module, \texttt{policy\_contract.py}, authoritatively defines what
the network sees: eight scalars, each clamped by a shared normaliser
(Table~\ref{tab:obs}).

\begin{table}[htbp]
\caption{The eight-field observation contract}
\label{tab:obs}
\centering\scriptsize
\setlength{\tabcolsep}{3pt}
\begin{tabular}{@{}L{2.25cm}L{1.4cm}L{2.0cm}L{1.45cm}@{}}
\toprule
Field & Meaning & Field & Meaning \\
\midrule
\texttt{lane\_error\_m}     & offset from lane centre & \texttt{front\_obst\_m} & nearest car ahead \\
\texttt{heading\_err\_deg}  & car vs.\ road angle & \texttt{red\_light}     & binary: light governs \\
\texttt{speed\_kmh}         & current road speed      & \texttt{curvature}      & bend $\approx$8\,m ahead \\
\texttt{target\_spd\_kmh}   & speed controller wants  & \texttt{stopped}        & binary: $<$1\,km/h \\
\bottomrule
\end{tabular}
\end{table}

Route-relative rather than world coordinates is what allows a policy
trained on one town to be evaluated on another. One gap is open: the
route-reference waypoint used for heading and curvature at training time
is not yet supplied identically at inference, so a train/serve skew
remains possible on two fields. Any result reported before it is closed
would be unsound.

\subsection{Action protocol and safety supervision}

The policy emits three values; the vehicle interface accepts seven.
Steering lies in $[-1,+1]$ and maps to $\pm70^{\circ}$; accelerator and
brake are clamped to $[0,1]$, and the runtime refuses to apply a
meaningful accelerator and brake simultaneously. The final layer is
linear --- all clamping is done by the runtime, not the network. The
remaining four fields (clutch, handbrake, wipers, gear) come from a
deterministic six-speed transmission model, and engine speed, which
CARLA does not simulate, is derived from road speed and gear ratio and
labelled \emph{estimated telemetry} wherever displayed.

A safety supervisor runs on the simulator host, outside the policy, and
forces a full brake on any of: collision, sensor timeout, a command
older than 0.5\,s, lane error above 2\,m, an obstacle within 7\,m, speed
above 50\,km/h, a red light, or invalid telemetry. It is the only
component permitted to override the network.

\subsection{Demonstration collection}

Two experts are used: CARLA's rule-based \texttt{BasicAgent}, and a
human at a six-zone cockpit. A route runs from spawn to a destination at
least 150\,m away, and every tick writes one JSONL row holding the eight
observations, the expert's three controls, telemetry and metadata. A
route is promoted only if it is collision-free, completed, and holds
$\geq$100 samples; crashes and interrupted runs are discarded rather
than truncated. \textbf{This rule is the mechanism by which the
rendering failures of Section~VI convert into missing data.}

All accepted recordings use Town05, with Town04 held out as an unseen
map so a good score cannot be explained by memorised road geometry.
Twelve Town04 teacher episodes recorded before this rule was adopted
were quarantined rather than deleted.

\subsection{Training pipeline}

Training is offline, in a Colab notebook, and produces a JSON artifact
the local driver loads directly. The model is a linear
$8\!\rightarrow\!24$ layer, $\tanh$, $24\!\rightarrow\!16$, $\tanh$, and
a linear $16\!\rightarrow\!3$ head: 667 parameters. Optimisation uses
AdamW ($\mathrm{lr}=10^{-3}$, weight decay $10^{-5}$) against mean
squared error on the three control channels, batch size 128, at most 80
epochs, early stopping after 12 without improvement. A single seed~(7)
is set across Python, NumPy and PyTorch; training runs on CPU.

The split is taken at route level, 70/15/15, never at frame level.
Consecutive frames within a route are near-identical, so a random frame
split places duplicates in both training and test sets and renders the
reported error meaningless. Each episode goes to exactly one split, and
an assertion confirms the sets are disjoint. The notebook generates its
normaliser verbatim from the contract file, so training-time
normalisation cannot diverge from inference-time.

\subsection{Deployability gate}

An exported artifact carries a policy-contract version, a
\texttt{deployable} flag and the episode names in each split. The live
driver refuses any artifact whose contract version is not v2 or whose
flag is false. At present only a synthetic smoke-test artifact exists
(contract v1, \texttt{deployable: false}), so it is refused and route
mode is unavailable --- and there is no scripted-agent fallback, so the
car is simply not driven autonomously. Training is hard-disabled at
runtime, so the deployed system cannot modify its weights.

One runtime behaviour must be declared, since it is not learned: cloning
over-represents braking, so the policy stalls on a clear road. A bounded
recovery capped at 0.35 accelerator releases it, disabled instantly by a
red light, an obstacle or invalid telemetry, and never edits the
weights. Any result must state whether it was active.

% =====================================================================
\section{Experimental Setup (partial)}

The scenario catalogue defines four categories so a policy can be
stressed across road type and traffic density without hard-coding a map:
urban (Town03/04/05/10HD, 30 NPCs), highway (Town05/06/07, 20), rural
(Town01/02, 12) and intersection (Town03/04, 40), with eight weather
presets and three pothole modes. These are \emph{configured} definitions
only --- \textbf{no category has yet completed a verified evaluation
run}, for the reasons in Section~VI.

\TODO{TODO --- SECTION V. Missing: a hardware table for both machines;
control-loop latency over the link (mean, p99); headless tick rate; the
commit hash for the reported runs.}

% =====================================================================
\section{The Rendering Bottleneck}

This section is the honest centre of the paper. The dominant obstacle
has not been the learning algorithm; it has been getting the environment
and the final camera image to render at all, for long enough to complete
a single route.

\subsection{The final image will not render at speed}

The visual cockpit composites four camera feeds alongside streamed
scenery. On an Intel Arc~140V integrated GPU this saturates the device:
the frame rate collapses as road speed rises, because faster travel
forces Unreal to stream new scenery while the same device is already
compositing the camera panel. The failure is not graceful: the image
stalls, the client falls behind the server tick, commands go stale.

Two mitigations are in place, both costing realism. The default neural
path runs \emph{headless}, with no camera feeds rendered at all. On the
visual build, throttle is removed above 42\,km/h and top speed capped at
48\,km/h purely to protect the GPU. The consequence must be stated: the
speed range for visual demonstrations is bounded by a rendering limit
rather than a driving decision, so no highway-speed behaviour can
currently be demonstrated.

\subsection{Environment rendering stalls crash the stack}

Environment rendering fails in three ways, all ending the episode:

\begin{itemize}
\setlength{\itemsep}{1pt}\setlength{\topsep}{2pt}
\setlength{\parskip}{0pt}\setlength{\partopsep}{0pt}
\item \textbf{Thermal termination.} The integrated GPU shares a thermal
package with the CPU. Sustained rendering drove a fatal hardware error
and a critical-thermal hibernation in one afternoon. No live temperature
counter is exposed here, so a health guard reads the Windows Event Log
instead and refuses to launch inside a 30- or 60-minute cooldown.
\item \textbf{Silent process death.} A rendering component wedges,
Windows kills it as a hung application, and the stack goes down. Nothing
was logged originally, because the launcher fired everything into
detached tasks whose output went nowhere. One supervisor now owns both
processes, writes real logs, probes CARLA's RPC port for liveness rather
than trusting a PID, and restarts with 5\,s\,$\rightarrow$\,120\,s
backoff.
\item \textbf{Orphaned renderers.} \texttt{CarlaUE4.exe} is a shim
around another executable, and the virtualenv \texttt{python.exe}
re-executes the real interpreter as a child. Killing the visible process
left the real renderer holding port 2000, so the next launch fought the
corpse of the previous one. The supervisor now kills process trees and
sweeps orphans at startup.
\end{itemize}

\subsection{Why this empties the training corpus}

These are not merely operational annoyances; they are why Section~VII is
blank. A rendering stall interrupts a route in progress, and by the rule
of Section~IV-C an interrupted route is discarded rather than kept --- a
deliberate choice, since a route ending in a crash teaches the policy
that crashing is normal. The two rules compose badly: the more the
environment fails to render, the fewer routes survive, and at present
\textbf{zero routes have been accepted} of twelve. The model is not
failing to learn; it has not been given anything to learn from.

\TODO{TODO --- SECTION VI. Qualitative; make quantitative before
submission. Needed: (a) a figure of FPS vs.\ road speed, headless vs.\
visual; (b) mean route survival time before a rendering-induced
termination; (c) counts per failure class; (d) package temperature at
termination (unobtainable --- state the limitation); (e) whether the
two-machine split relieves the bottleneck.}

% =====================================================================
\section{Results}

\UNWRITTEN{This section is not yet written.}

The corpus does not exist, for the reasons in Section~VI, and a training
loss against the smoke-test artifact would be misleading. To prevent
metric selection after the fact, the reporting format is fixed now:
collisions per 100\,km, lane departures per km, route completion rate,
safety interventions per route, and per-channel MAE --- each for the
rule-based baseline and the learned policy on the held-out map, and
every one of them \emph{pending}.

\TODO{TODO --- SECTION VII. Prerequisites, in order: (1) stabilise
rendering so a 150\,m route survives end-to-end; (2) accept $\geq$12
Town05 routes; (3) close the train--serve skew of Section~IV-A; (4)
export a contract-v2 artifact with \texttt{deployable = true}; (5) run
the Town04 evaluation with the stall assist on and off. No prose before
(5) completes.}

% =====================================================================
\section{Discussion}

\UNWRITTEN{Stub. Only observations independent of quantitative results
are recorded.}

The project's time went to rendering and process management rather than
modelling, worth reporting because it is absent from the work that
motivated ours. Every result in \cite{lbc, transfuser} presupposes a
simulator that renders reliably; where that fails, the bottleneck moves
somewhere the literature does not discuss.

\TODO{TODO --- SECTION VIII. Once VII exists: (a) a leakage demo ---
error under route- vs.\ frame-level splitting; (b) a failure taxonomy
from Town04; (c) what eight scalars cannot represent (pedestrians,
signs, any cue needing vision); (d) whether headless collection biases
the policy relative to the visual build.}

% =====================================================================
\section{Conclusion and Future Work}

\UNWRITTEN{[Not written. To be drafted after Sections~VI and~VII are
complete.]}

Placeholder for planned direction: resolving the rendering bottleneck,
if necessary by relocating the renderer to dedicated hardware; migration
from eight scalars to RGB, depth and semantic segmentation; conditional
imitation learning~\cite{cil}; DAgger-style correction~\cite{dagger};
and a curriculum across the catalogue. Reward logging already runs each
tick, but no online loop consumes it.

% =====================================================================
\begin{thebibliography}{99}
\footnotesize
\setlength{\itemsep}{0pt}\setlength{\parsep}{0pt}

\bibitem{carla} A. Dosovitskiy \textit{et al.}, ``CARLA: An Open Urban
Driving Simulator,'' \emph{CoRL}, 2017. arXiv:1711.03938.

\bibitem{bojarski} M. Bojarski \textit{et al.}, ``End to End Learning
for Self-Driving Cars,'' NVIDIA, 2016. arXiv:1604.07316.

\bibitem{cil} F. Codevilla \textit{et al.}, ``End-to-end Driving via
Conditional Imitation Learning,'' \emph{ICRA}, 2018. arXiv:1710.02410.

\bibitem{limits} F. Codevilla \textit{et al.}, ``Exploring the
Limitations of Behavior Cloning for Autonomous Driving,'' \emph{ICCV},
2019.

\bibitem{dagger} S. Ross, G. Gordon, and J. A. Bagnell, ``A Reduction of
Imitation Learning and Structured Prediction to No-Regret Online
Learning,'' \emph{AISTATS}, 2011. arXiv:1011.0686.

\bibitem{lbc} D. Chen \textit{et al.}, ``Learning by Cheating,''
\emph{CoRL}, 2019. arXiv:1912.12294.

\bibitem{transfuser} A. Prakash, K. Chitta, and A. Geiger,
``Multi-Modal Fusion Transformer for End-to-End Autonomous Driving,''
\emph{CVPR}, 2021.

\end{thebibliography}

\TODOflow{TODO --- REFERENCES incomplete: add CARLA docs, OpenDRIVE,
rendering refs.}

\end{document}
